\section{Methodology}\label{sec:method}

In this section I present the benchmark dataset and adversarial algorithm utilized to evaluate OOD and adversarial robustness of a (pruned) network.

\subsection{CIFAR Corrupted Benchmark}

The CIFAR-10 and CIFAR-100 datasets have been used as benchmarks for decades in the field of computer vision \cite{cifar10}. The two datasets consits of
60,000 32x32 RGB images with 10 and 100 different classes, respectively. The images are uniformly distributed across the classes and objects are at the center, thus
it serves as a clean baseline without any modification that may happen in real-world scenarios.

To facilitate the prediction of the impact of common image corruptions the trained model, the CIFAR-10-C and CIFAR-100-C datasets were
 constructed \cite{cifar10c}. The authors applied 19 different types of corruptions at varying intensities on the test sets of CIFAR-10 and CIFAR-100, respectively. 
 More specifically, corruptions are categories into followig categories: noise, blur, weather conditions, and digital corruptions, while intensity levels span from 1 to 5.
 Figure \ref{fig:cifar10c} shows an example for every category, applied on the same image from CIFAR-10 dataset.


 \begin{figure}[H]
    \centering
    \includegraphics[width=1\linewidth]{figs/cifar10c_grid_idx100_sev5.png}
    \caption{Different types of corruptions with severity level 5 on a CIFAR-10 32x32 RGB image}
    \label{fig:cifar10c}
\end{figure}\

In the context of the OOD generalization problem mentioned in section \ref{subsec:ood&adv}, these real-world distortions cause
a covariate shift in the distribution of the test data. In other words, only input features suffer a distributional shift, while
the underlying concept - the class of photographed object in this case - remains the same.
\subsection{Fast Adaptive Boundary}

\subsection{Fourier Sensitivity Method}